%==============================================================================
% GMAO Chatbot - Architecture LLM Text-to-SQL
% Documentation Technique pour Superviseur
% Date: 27 Janvier 2026
%==============================================================================

\documentclass[12pt,a4paper]{article}

%------------------------------------------------------------------------------
% PACKAGES
%------------------------------------------------------------------------------
\usepackage[utf8]{inputenc}
\usepackage[T1]{fontenc}
\usepackage[french]{babel}
\usepackage{geometry}
\usepackage{graphicx}
\usepackage{xcolor}
\usepackage{tikz}
\usepackage{tcolorbox}
\usepackage{listings}
\usepackage{booktabs}
\usepackage{array}
\usepackage{tabularx}
\usepackage{colortbl}
\usepackage{hyperref}
\usepackage{fontawesome5}
\usepackage{enumitem}
\usepackage{fancyhdr}
\usepackage{titlesec}
\usepackage{float}
\usepackage{caption}
\usepackage{subcaption}
\usepackage{amsmath}
\usepackage{amssymb}

%------------------------------------------------------------------------------
% CONFIGURATION DE LA PAGE
%------------------------------------------------------------------------------
\geometry{
    left=2cm,
    right=2cm,
    top=2.5cm,
    bottom=2.5cm
}

%------------------------------------------------------------------------------
% COULEURS PERSONNALISÉES
%------------------------------------------------------------------------------
% Couleurs principales
\definecolor{primaryblue}{RGB}{41, 98, 255}
\definecolor{secondarygreen}{RGB}{46, 204, 113}
\definecolor{accentorange}{RGB}{243, 156, 18}
\definecolor{warningred}{RGB}{231, 76, 60}
\definecolor{darkgray}{RGB}{52, 73, 94}
\definecolor{lightgray}{RGB}{236, 240, 241}
\definecolor{codebg}{RGB}{45, 45, 45}
\definecolor{codegreen}{RGB}{152, 195, 121}
\definecolor{codepurple}{RGB}{198, 120, 221}
\definecolor{codeyellow}{RGB}{229, 192, 123}
\definecolor{codeblue}{RGB}{97, 175, 239}

% Couleurs pour les boîtes
\definecolor{infobox}{RGB}{52, 152, 219}
\definecolor{successbox}{RGB}{39, 174, 96}
\definecolor{warningbox}{RGB}{241, 196, 15}
\definecolor{dangerbox}{RGB}{192, 57, 43}

%------------------------------------------------------------------------------
% CONFIGURATION DES LIENS
%------------------------------------------------------------------------------
\hypersetup{
    colorlinks=true,
    linkcolor=primaryblue,
    urlcolor=primaryblue,
    citecolor=primaryblue
}

%------------------------------------------------------------------------------
% STYLE DES TITRES
%------------------------------------------------------------------------------
\titleformat{\section}
    {\color{primaryblue}\LARGE\bfseries}
    {\thesection}{1em}{}[\titlerule]

\titleformat{\subsection}
    {\color{darkgray}\Large\bfseries}
    {\thesubsection}{1em}{}

\titleformat{\subsubsection}
    {\color{darkgray}\large\bfseries}
    {\thesubsubsection}{1em}{}

%------------------------------------------------------------------------------
% BOÎTES COLORÉES (tcolorbox)
%------------------------------------------------------------------------------
\tcbuselibrary{skins,breakable}

% Boîte Info
\newtcolorbox{infobox}[1][]{
    enhanced,
    colback=infobox!10,
    colframe=infobox,
    fonttitle=\bfseries,
    title={\faInfoCircle\ Information},
    #1
}

% Boîte Succès
\newtcolorbox{successbox}[1][]{
    enhanced,
    colback=successbox!10,
    colframe=successbox,
    fonttitle=\bfseries,
    title={\faCheckCircle\ Avantage},
    #1
}

% Boîte Avertissement
\newtcolorbox{warningbox}[1][]{
    enhanced,
    colback=warningbox!10,
    colframe=warningbox,
    fonttitle=\bfseries,
    title={\faExclamationTriangle\ Attention},
    #1
}

% Boîte Danger
\newtcolorbox{dangerbox}[1][]{
    enhanced,
    colback=dangerbox!10,
    colframe=dangerbox,
    fonttitle=\bfseries,
    title={\faTimesCircle\ Interdit},
    #1
}

% Boîte Code/Technique
\newtcolorbox{techbox}[1][]{
    enhanced,
    colback=codebg,
    colframe=darkgray,
    fonttitle=\bfseries\color{white},
    coltitle=white,
    title={\faCode\ Code},
    #1
}

% Boîte Définition
\newtcolorbox{defbox}[2][]{
    enhanced,
    colback=lightgray,
    colframe=primaryblue,
    fonttitle=\bfseries,
    title={\faBook\ #2},
    #1
}

%------------------------------------------------------------------------------
% CONFIGURATION DES LISTINGS (CODE)
%------------------------------------------------------------------------------
\lstset{
    backgroundcolor=\color{codebg},
    basicstyle=\ttfamily\small\color{white},
    keywordstyle=\color{codepurple}\bfseries,
    stringstyle=\color{codegreen},
    commentstyle=\color{gray}\itshape,
    numberstyle=\tiny\color{gray},
    numbers=left,
    numbersep=8pt,
    breaklines=true,
    breakatwhitespace=true,
    tabsize=4,
    showstringspaces=false,
    frame=single,
    rulecolor=\color{darkgray},
    captionpos=b,
    xleftmargin=15pt,
    framexleftmargin=15pt
}

\lstdefinelanguage{SQL}{
    keywords={SELECT, FROM, WHERE, JOIN, LEFT, RIGHT, INNER, ON, AND, OR, NOT, IN, AS, TOP, ORDER, BY, GROUP, HAVING, COUNT, SUM, AVG, MAX, MIN, CREATE, TABLE, INSERT, UPDATE, DELETE, DROP, TRUNCATE, ALTER, PRIMARY, KEY, FOREIGN, REFERENCES, INT, NVARCHAR, BIT, DATETIME, DEFAULT, NULL, CAST, GETDATE, DATE, DATEADD},
    keywordstyle=\color{codeblue}\bfseries,
    sensitive=false,
    morestring=[b]',
    morecomment=[l]{--},
}

\lstdefinelanguage{Python}{
    keywords={def, return, if, else, elif, for, while, in, import, from, as, class, try, except, with, True, False, None, and, or, not},
    keywordstyle=\color{codepurple}\bfseries,
    sensitive=true,
    morestring=[b]",
    morestring=[b]',
    morecomment=[l]{\#},
}

\lstdefinelanguage{JSON}{
    morestring=[b]",
    stringstyle=\color{codegreen},
    literate=
     *{:}{{{\color{codeyellow}:}}}{1}
      {,}{{{\color{white},}}}{1}
      {\{}{{{\color{codepurple}\{}}}{1}
      {\}}{{{\color{codepurple}\}}}}{1}
      {[}{{{\color{codepurple}[}}}{1}
      {]}{{{\color{codepurple}]}}}{1},
}

%------------------------------------------------------------------------------
% EN-TÊTE ET PIED DE PAGE
%------------------------------------------------------------------------------
\pagestyle{fancy}
\fancyhf{}
\fancyhead[L]{\textcolor{darkgray}{\small GMAO Chatbot - Architecture}}
\fancyhead[R]{\textcolor{darkgray}{\small \thepage}}
\fancyfoot[C]{\textcolor{gray}{\small Document Technique - Janvier 2026}}
\renewcommand{\headrulewidth}{0.4pt}
\renewcommand{\footrulewidth}{0.4pt}

%------------------------------------------------------------------------------
% TikZ LIBRARIES
%------------------------------------------------------------------------------
\usetikzlibrary{shapes.geometric, arrows.meta, positioning, calc, shadows, fit, backgrounds}

%==============================================================================
% DÉBUT DU DOCUMENT
%==============================================================================
\begin{document}

%------------------------------------------------------------------------------
% PAGE DE TITRE
%------------------------------------------------------------------------------
\begin{titlepage}
    \centering
    \vspace*{2cm}
    
    % Logo/Icône
    {\Huge\color{primaryblue}\faRobot}
    
    \vspace{1cm}
    
    {\Huge\bfseries\color{primaryblue} GMAO Chatbot}
    
    \vspace{0.5cm}
    
    {\LARGE\color{darkgray} Architecture LLM Text-to-SQL}
    
    \vspace{2cm}
    
    \begin{tikzpicture}
        \node[draw=primaryblue, fill=primaryblue!10, rounded corners=10pt, 
              minimum width=12cm, minimum height=3cm, align=center] {
            {\Large\bfseries Documentation Technique}\\[0.3cm]
            {\large Assistant IA pour la Gestion de Maintenance}
        };
    \end{tikzpicture}
    
    \vspace{2cm}
    
    \begin{tabular}{ll}
        \faCalendar & \textbf{Date :} 27 Janvier 2026 \\[0.3cm]
        \faCode & \textbf{Version :} 1.0 \\[0.3cm]
        \faServer & \textbf{Environnement :} Développement \\[0.3cm]
        \faBrain & \textbf{Modèle IA :} Gemma 3 27B (OpenRouter)
    \end{tabular}
    
    \vfill
    
    {\large\color{gray} Projet de Stage - GMAO}
    
\end{titlepage}

%------------------------------------------------------------------------------
% TABLE DES MATIÈRES
%------------------------------------------------------------------------------
\tableofcontents
\newpage

%==============================================================================
% SECTION 1 : VUE D'ENSEMBLE
%==============================================================================
\section{\faEye\ Vue d'Ensemble}

\subsection{Qu'est-ce que le GMAO Chatbot ?}

\begin{defbox}{Définition}
Le \textbf{GMAO Chatbot} est un assistant intelligent basé sur l'Intelligence Artificielle (IA) qui permet aux superviseurs et techniciens d'interroger la base de données GMAO en \textbf{langage naturel français}.
\end{defbox}

\begin{successbox}
\textbf{Exemple d'utilisation :}\\
L'utilisateur demande : \textit{"Quelles sont les interventions en cours ?"}\\
Le chatbot répond automatiquement avec la liste des interventions actives.
\end{successbox}

\subsection{Approche Technique : Text-to-SQL}

Notre chatbot utilise l'approche \textbf{Text-to-SQL}, qui convertit les questions en langage naturel en requêtes SQL exécutables sur la base de données.

\begin{figure}[H]
\centering
\begin{tikzpicture}[
    node distance=1.2cm,
    box/.style={rectangle, rounded corners, draw=#1, fill=#1!20, 
                minimum width=10cm, minimum height=1cm, align=center,
                font=\small},
    arrow/.style={-{Stealth[length=3mm]}, thick, color=darkgray}
]

% Nodes
\node[box=primaryblue] (q1) {\faUser\ \textbf{Question Utilisateur (Français)}\\
"Quelles sont les interventions en cours ?"};

\node[box=accentorange, below=of q1] (llm) {\faBrain\ \textbf{LLM (Large Language Model)}\\
Analyse et convertit en SQL};

\node[box=secondarygreen, below=of llm] (sql) {\faDatabase\ \textbf{Requête SQL Générée}\\
\texttt{SELECT * FROM DemandeInterventions WHERE StatusId = 2}};

\node[box=darkgray, below=of sql] (db) {\faServer\ \textbf{Exécution sur SQL Server}\\
Récupération des données};

\node[box=primaryblue, below=of db] (resp) {\faComment\ \textbf{Réponse Formatée}\\
"Il y a 3 interventions en cours : ..."};

% Arrows
\draw[arrow] (q1) -- (llm);
\draw[arrow] (llm) -- (sql);
\draw[arrow] (sql) -- (db);
\draw[arrow] (db) -- (resp);

\end{tikzpicture}
\caption{Flux de traitement Text-to-SQL}
\end{figure}

\newpage

%==============================================================================
% SECTION 2 : ARCHITECTURE IA - TEXT-TO-SQL vs RAG
%==============================================================================
\section{\faBrain\ Architecture IA : Pourquoi Text-to-SQL ?}

\subsection{Deux Approches Possibles}

Il existe deux approches principales pour créer un chatbot IA connecté à une base de données :

\begin{figure}[H]
\centering
\begin{tikzpicture}[
    node distance=0.5cm,
    techbox/.style={rectangle, rounded corners=5pt, draw=#1, fill=#1!15,
                    minimum width=6cm, minimum height=0.8cm, align=center, font=\small},
    label/.style={font=\bfseries\large, color=#1}
]

% RAG Section
\node[label=warningred] (ragtitle) at (-4, 4) {\faTimesCircle\ RAG (Non utilisé)};

\node[techbox=warningred, below=of ragtitle] (rag1) {Documents $\rightarrow$ Chunking};
\node[techbox=warningred, below=of rag1] (rag2) {Embeddings $\rightarrow$ Vector DB};
\node[techbox=warningred, below=of rag2] (rag3) {Recherche Similarité};
\node[techbox=warningred, below=of rag3] (rag4) {LLM $\rightarrow$ Réponse};

% Text-to-SQL Section
\node[label=successbox, right=3cm of ragtitle] (sqltitle) {\faCheckCircle\ Text-to-SQL (Utilisé)};

\node[techbox=successbox, below=of sqltitle] (sql1) {Question Naturelle};
\node[techbox=successbox, below=of sql1] (sql2) {LLM + Schéma DB};
\node[techbox=successbox, below=of sql2] (sql3) {Requête SQL};
\node[techbox=successbox, below=of sql3] (sql4) {Résultat Exact};

% Arrows
\foreach \i/\j in {rag1/rag2, rag2/rag3, rag3/rag4} {
    \draw[-{Stealth}, thick, warningred] (\i) -- (\j);
}
\foreach \i/\j in {sql1/sql2, sql2/sql3, sql3/sql4} {
    \draw[-{Stealth}, thick, successbox] (\i) -- (\j);
}

\end{tikzpicture}
\caption{Comparaison RAG vs Text-to-SQL}
\end{figure}

\subsection{Qu'est-ce que le RAG ?}

\begin{defbox}{RAG - Retrieval-Augmented Generation}
Le RAG est une technique qui combine :
\begin{enumerate}
    \item Une \textbf{base de données vectorielle} (Pinecone, Chroma, Weaviate)
    \item Des \textbf{embeddings} (représentations numériques du texte)
    \item Un \textbf{LLM} pour générer des réponses
\end{enumerate}
\textbf{Idéal pour :} Documents textuels, FAQ, manuels, recherche sémantique.
\end{defbox}

\subsection{Pourquoi nous n'utilisons PAS le RAG ?}

\begin{dangerbox}[title={\faTimesCircle\ Problèmes du RAG pour GMAO}]
\begin{itemize}[leftmargin=*]
    \item \textbf{Données structurées :} RAG est conçu pour du texte, pas pour des tableaux SQL
    \item \textbf{Synchronisation :} Chaque modification en DB nécessite une re-indexation
    \item \textbf{Précision :} "Combien d'interventions ?" $\rightarrow$ RAG donne une approximation
    \item \textbf{Requêtes complexes :} RAG ne peut pas faire GROUP BY, SUM, COUNT, JOIN
    \item \textbf{Coût élevé :} Vector DB + Embeddings $\approx$ \textbf{140€/mois}
\end{itemize}
\end{dangerbox}

\subsection{Pourquoi Text-to-SQL est Idéal ?}

\begin{successbox}[title={\faCheckCircle\ Avantages du Text-to-SQL pour GMAO}]
\begin{itemize}[leftmargin=*]
    \item \textbf{Temps réel :} Interroge directement la base de données
    \item \textbf{Précision 100\% :} SQL retourne les chiffres exacts
    \item \textbf{Requêtes complexes :} Agrégations, jointures, filtres natifs
    \item \textbf{Pas de Vector DB :} Économie d'infrastructure
    \item \textbf{Coût faible :} Seulement LLM $\approx$ \textbf{15-20€/mois}
\end{itemize}
\end{successbox}

\subsection{Tableau Comparatif Détaillé}

\begin{table}[H]
\centering
\rowcolors{2}{lightgray}{white}
\begin{tabular}{|>{\bfseries}l|c|c|}
\hline
\rowcolor{primaryblue}
\textcolor{white}{\textbf{Critère}} & \textcolor{white}{\textbf{Text-to-SQL} \faCheckCircle} & \textcolor{white}{\textbf{RAG} \faTimesCircle} \\
\hline
Données en temps réel & \cellcolor{successbox!30}\faCheck\ Toujours à jour & \cellcolor{dangerbox!30}\faTimes\ Synchronisation requise \\
\hline
Précision des chiffres & \cellcolor{successbox!30}\faCheck\ 100\% exacte & \cellcolor{warningbox!30}\faExclamation\ Approximative \\
\hline
Complexité d'installation & \cellcolor{successbox!30}\faCheck\ Simple & \cellcolor{dangerbox!30}\faTimes\ Complexe \\
\hline
Coût d'infrastructure & \cellcolor{successbox!30}\faCheck\ $\sim$20€/mois & \cellcolor{dangerbox!30}\faTimes\ $\sim$140€/mois \\
\hline
Requêtes analytiques & \cellcolor{successbox!30}\faCheck\ Excellent & \cellcolor{dangerbox!30}\faTimes\ Limité \\
\hline
\end{tabular}
\caption{Comparaison Text-to-SQL vs RAG}
\end{table}

\newpage

%==============================================================================
% SECTION 3 : ARCHITECTURE TECHNIQUE COMPLÈTE
%==============================================================================
\section{\faCogs\ Architecture Technique Complète}

\subsection{Vue d'Ensemble du Système}

\begin{figure}[H]
\centering
\begin{tikzpicture}[
    scale=0.85,
    transform shape,
    node distance=0.8cm,
    component/.style={rectangle, rounded corners=8pt, draw=#1, fill=#1!20,
                      minimum width=11cm, minimum height=1.5cm, align=center,
                      font=\small, drop shadow},
    subcomp/.style={rectangle, rounded corners=3pt, draw=gray, fill=white,
                    minimum width=3cm, minimum height=0.6cm, align=center,
                    font=\footnotesize},
    arrow/.style={-{Stealth[length=3mm]}, thick, color=darkgray},
    label/.style={font=\bfseries\footnotesize, color=white}
]

% Frontend
\node[component=primaryblue] (frontend) {
    \textcolor{primaryblue}{\faDesktop\ \textbf{FRONTEND - Angular 17+}}\\
    http://localhost:4200
};

% API .NET
\node[component=secondarygreen, below=1.5cm of frontend] (api) {
    \textcolor{secondarygreen}{\faServer\ \textbf{API .NET Core 9.0}}\\
    http://localhost:5000
};

% Python Service
\node[component=accentorange, below=1.5cm of api] (python) {
    \textcolor{accentorange}{\faPython\ \textbf{PYTHON CHATBOT SERVICE - FastAPI}}\\
    http://localhost:5002
};

% Database
\node[component=darkgray, below=1.5cm of python] (db) {
    \textcolor{darkgray}{\faDatabase\ \textbf{SQL SERVER DATABASE}}\\
    localhost$\backslash$SQLEXPRESS - Database: gmao
};

% LLM Provider
\node[component=codepurple, right=0.5cm of python, minimum width=4cm] (llm) {
    \textcolor{codepurple}{\faBrain\ \textbf{LLM}}\\
    OpenRouter\\
    Gemma 3 27B
};

% Arrows
\draw[arrow] (frontend) -- node[right, font=\footnotesize] {HTTP POST /query} (api);
\draw[arrow] (api) -- node[right, font=\footnotesize] {HTTP POST /query} (python);
\draw[arrow] (python) -- node[right, font=\footnotesize] {SQL Queries} (db);
\draw[arrow] (python) -- node[above, font=\footnotesize] {API Calls} (llm);

% Labels on left
\node[left=0.3cm of frontend, font=\footnotesize\color{gray}] {Interface Utilisateur};
\node[left=0.3cm of api, font=\footnotesize\color{gray}] {Authentification JWT};
\node[left=0.3cm of python, font=\footnotesize\color{gray}] {Traitement NLP + SQL};
\node[left=0.3cm of db, font=\footnotesize\color{gray}] {Stockage des données};

\end{tikzpicture}
\caption{Architecture Globale du Système GMAO Chatbot}
\end{figure}

\subsection{Composants du Service Python (Cœur du Chatbot)}

\begin{figure}[H]
\centering
\begin{tikzpicture}[
    node distance=0.6cm,
    step/.style={rectangle, rounded corners=5pt, draw=accentorange, fill=accentorange!15,
                 minimum width=10cm, minimum height=1cm, align=center, font=\small},
    num/.style={circle, draw=accentorange, fill=accentorange, text=white, 
                minimum size=0.6cm, font=\bfseries\small}
]

\node[step] (s1) {\textbf{Query Preprocessor} - Normalisation, détection langue, dates françaises};
\node[num, left=0.3cm of s1] {1};

\node[step, below=of s1] (s2) {\textbf{Predefined Query Matcher} - Patterns regex pour requêtes communes};
\node[num, left=0.3cm of s2] {2};

\node[step, below=of s2] (s3) {\textbf{LLM Text-to-SQL Engine} - Génération SQL par IA};
\node[num, left=0.3cm of s3] {3};

\node[step, below=of s3] (s4) {\textbf{SQL Validator} - Validation sécurité (SELECT only)};
\node[num, left=0.3cm of s4] {4};

\node[step, below=of s4] (s5) {\textbf{SQL Executor} - Exécution sur SQL Server};
\node[num, left=0.3cm of s5] {5};

\node[step, below=of s5] (s6) {\textbf{Response Formatter} - Formatage réponse en français};
\node[num, left=0.3cm of s6] {6};

% Arrows
\foreach \i/\j in {s1/s2, s2/s3, s3/s4, s4/s5, s5/s6} {
    \draw[-{Stealth}, thick, accentorange] (\i) -- (\j);
}

\end{tikzpicture}
\caption{Pipeline de Traitement du Service Python}
\end{figure}

\newpage

%==============================================================================
% SECTION 4 : ARCHITECTURE HYBRIDE (3 NIVEAUX)
%==============================================================================
\section{\faLayerGroup\ Architecture Hybride : 3 Niveaux d'Optimisation}

Notre chatbot utilise une \textbf{approche hybride} pour optimiser les performances et réduire les coûts.

\begin{figure}[H]
\centering
\begin{tikzpicture}[
    node distance=0.8cm,
    level/.style={rectangle, rounded corners=10pt, draw=#1, fill=#1!20,
                  minimum width=12cm, minimum height=2cm, align=center,
                  font=\small, drop shadow},
    arrow/.style={-{Stealth[length=3mm]}, thick},
    label/.style={font=\bfseries}
]

% Input
\node[rectangle, rounded corners=5pt, draw=primaryblue, fill=primaryblue!30,
      minimum width=8cm, minimum height=0.8cm] (input) {
    \faUser\ \textbf{Question Utilisateur}
};

% Level 1
\node[level=successbox, below=1cm of input] (l1) {
    \textbf{\faRocket\ NIVEAU 1 : Pattern Matching (Regex)}\\[0.2cm]
    Questions fréquentes prédéfinies $\bullet$ Réponse $\sim$100ms $\bullet$ \textbf{Coût : 0€}
};

% Level 2
\node[level=infobox, below=of l1] (l2) {
    \textbf{\faUserTag\ NIVEAU 2 : Détection d'Entités}\\[0.2cm]
    "Interventions de Mohamed" $\bullet$ SQL paramétré $\bullet$ \textbf{Coût : 0€}
};

% Level 3
\node[level=accentorange, below=of l2] (l3) {
    \textbf{\faBrain\ NIVEAU 3 : LLM Text-to-SQL}\\[0.2cm]
    Questions complexes/nouvelles $\bullet$ Génération SQL $\bullet$ \textbf{Coût : $\sim$0.001€/requête}
};

% Output
\node[rectangle, rounded corners=5pt, draw=secondarygreen, fill=secondarygreen!30,
      minimum width=8cm, minimum height=0.8cm, below=1cm of l3] (output) {
    \faComment\ \textbf{Réponse Formatée}
};

% Arrows
\draw[arrow, primaryblue] (input) -- (l1);
\draw[arrow, gray] (l1) -- node[right, font=\footnotesize] {si pas de match} (l2);
\draw[arrow, gray] (l2) -- node[right, font=\footnotesize] {si pas de match} (l3);
\draw[arrow, successbox] (l3) -- (output);

% Direct arrows for matches
\draw[arrow, successbox, dashed] (l1.east) -- ++(2,0) |- (output.east);
\draw[arrow, infobox, dashed] (l2.east) -- ++(1.5,0) |- ($(output.east)-(0,0.2)$);

% Percentages
\node[right=2.5cm of l1, font=\bfseries\color{successbox}] {$\sim$50\%};
\node[right=2cm of l2, font=\bfseries\color{infobox}] {$\sim$20\%};
\node[right=0.5cm of l3, font=\bfseries\color{accentorange}] {$\sim$30\%};

\end{tikzpicture}
\caption{Architecture Hybride à 3 Niveaux}
\end{figure}

\begin{infobox}[title={\faChartPie\ Résultats de cette Architecture}]
\begin{tabular}{ll}
    \textbf{Requêtes sans LLM :} & $\sim$70\% (niveaux 1 + 2) \\
    \textbf{Économie :} & $\sim$70\% de réduction des coûts LLM \\
    \textbf{Temps de réponse moyen :} & $\sim$500ms (vs 5s full LLM) \\
    \textbf{Précision :} & 100\% pour SQL, $\sim$90\% pour LLM \\
\end{tabular}
\end{infobox}

\newpage

%==============================================================================
% SECTION 5 : WORKFLOW DÉTAILLÉ
%==============================================================================
\section{\faProjectDiagram\ Workflow Détaillé}

\subsection{Étape 1 : Réception de la Question}

\begin{lstlisting}[language=JSON, caption=Requête entrante]
{
    "query": "Quelles sont les interventions d'aujourd'hui ?"
}
\end{lstlisting}

\subsection{Étape 2 : Prétraitement}

\begin{lstlisting}[language=Python, caption=Prétraitement de la question]
# Normalisation
question = question.lower().strip()
# "quelles sont les interventions d'aujourd'hui ?"

# Detection des dates francaises
# "aujourd'hui" -> CAST(GETDATE() AS DATE)
# "hier" -> DATEADD(day, -1, CAST(GETDATE() AS DATE))
# "janvier 2026" -> '2026-01-01' AND '2026-01-31'
\end{lstlisting}

\subsection{Étape 3 : Matching des Patterns Prédéfinis}

\begin{lstlisting}[language=Python, caption=Patterns prédéfinis]
PREDEFINED_QUERIES = {
    "interventions": {
        "patterns": [r"intervention", r"demande", r"ticket"],
        "sql": """
            SELECT TOP 20
                d.Titre, s.Statut, d.DateCreation
            FROM DemandeInterventions d
            LEFT JOIN Statuses s ON d.StatusId = s.Id
            WHERE d.IsActif = 1
            ORDER BY d.DateCreation DESC
        """
    }
}
\end{lstlisting}

\subsection{Étape 4 : Génération SQL par LLM}

Si aucun pattern ne correspond, le LLM génère la requête :

\begin{lstlisting}[language=Python, caption=Prompt Engineering pour le LLM]
SYSTEM_PROMPT = f"""
Tu es un expert SQL Server pour une base de donnees GMAO.
Genere UNIQUEMENT une requete SQL SELECT valide.

## Schema de la base de donnees:
{DATABASE_SCHEMA}

## Regles:
1. Utilise uniquement SELECT (pas d'INSERT, UPDATE, DELETE)
2. Limite les resultats a TOP 20
3. Utilise les JOINs appropries

## Question utilisateur:
{user_question}
"""
\end{lstlisting}

\subsection{Étape 5 : Validation Sécurité}

\begin{dangerbox}[title={\faShieldAlt\ Validation SQL - Sécurité}]
\begin{lstlisting}[language=Python, caption=Validation de sécurité]
FORBIDDEN_KEYWORDS = [
    'INSERT', 'UPDATE', 'DELETE', 'DROP', 'TRUNCATE',
    'ALTER', 'CREATE', 'EXEC', 'EXECUTE', 'xp_', 'sp_',
    '--', ';', '/*', '*/'
]

def validate_sql(sql: str) -> bool:
    for keyword in FORBIDDEN_KEYWORDS:
        if keyword in sql.upper():
            return False  # BLOQUE!
    return sql.strip().upper().startswith('SELECT')
\end{lstlisting}
\end{dangerbox}

\subsection{Étape 6 : Réponse Finale}

\begin{lstlisting}[language=JSON, caption=Réponse du chatbot]
{
    "response": "Il y a 1 equipement en panne:\n
                 * Tour CNC vertical (EQ-007)",
    "metadata": {
        "ai_provider": "openrouter",
        "sql_executed": true,
        "results_count": 1,
        "processing_time_ms": 4339
    },
    "suggestions": [
        "Quelles sont les interventions d'aujourd'hui ?",
        "Combien d'equipements sont en panne ?"
    ]
}
\end{lstlisting}

\newpage

%==============================================================================
% SECTION 6 : TECHNOLOGIES UTILISÉES
%==============================================================================
\section{\faTools\ Technologies Utilisées}

\subsection{Stack Technique Complet}

\begin{table}[H]
\centering
\rowcolors{2}{lightgray}{white}
\begin{tabular}{|>{\bfseries}l|l|l|}
\hline
\rowcolor{primaryblue}
\textcolor{white}{\textbf{Composant}} & \textcolor{white}{\textbf{Technologie}} & \textcolor{white}{\textbf{Version}} \\
\hline
\multicolumn{3}{|c|}{\cellcolor{accentorange!30}\textbf{Backend Python (Chatbot)}} \\
\hline
Framework API & FastAPI & 0.104+ \\
Serveur & Uvicorn & 0.24+ \\
Connexion DB & PyODBC & 5.0+ \\
Validation & Pydantic & 2.0+ \\
HTTP Client & Requests & 2.31+ \\
\hline
\multicolumn{3}{|c|}{\cellcolor{secondarygreen!30}\textbf{Backend .NET (API)}} \\
\hline
Framework & .NET Core & 9.0 \\
ORM & Entity Framework Core & 9.0 \\
Auth & ASP.NET Identity + JWT & 9.0 \\
\hline
\multicolumn{3}{|c|}{\cellcolor{primaryblue!30}\textbf{Frontend}} \\
\hline
Framework & Angular & 17+ \\
Langage & TypeScript & 5.0+ \\
UI Components & Angular Material & 17+ \\
\hline
\multicolumn{3}{|c|}{\cellcolor{darkgray!30}\textbf{Base de Données}} \\
\hline
SGBD & SQL Server Express & 2019+ \\
Driver & ODBC Driver & 17+ \\
\hline
\multicolumn{3}{|c|}{\cellcolor{codepurple!30}\textbf{LLM Provider}} \\
\hline
Fournisseur & OpenRouter & Free Tier \\
Modèle & google/gemma-3-27b-it & Free \\
\hline
\end{tabular}
\caption{Stack Technique du GMAO Chatbot}
\end{table}

\newpage

%==============================================================================
% SECTION 7 : ÉTAT ACTUEL VS PRODUCTION
%==============================================================================
\section{\faChartLine\ État Actuel vs Production}

\subsection{Configuration Actuelle (Développement)}

\begin{table}[H]
\centering
\rowcolors{2}{lightgray}{white}
\begin{tabular}{|>{\bfseries}l|c|l|}
\hline
\rowcolor{accentorange}
\textcolor{white}{\textbf{Aspect}} & \textcolor{white}{\textbf{État}} & \textcolor{white}{\textbf{Note}} \\
\hline
LLM Provider & OpenRouter Free & Limité en requêtes/minute \\
Modèle & Gemma 3 27B & Gratuit mais moins précis \\
Temps de réponse & 4-17 secondes & Acceptable pour dev \\
Disponibilité & $\sim$99\% & Dépend d'OpenRouter \\
\rowcolor{successbox!30}
Coût & \textbf{0€/mois} & Gratuit \\
\hline
\end{tabular}
\caption{Configuration Actuelle}
\end{table}

\subsection{Limitations Actuelles}

\begin{warningbox}[title={\faExclamationTriangle\ Limitations à Résoudre}]
\begin{enumerate}
    \item \textbf{Rate Limits :} $\sim$10-20 requêtes/minute sur tier gratuit
    \item \textbf{Latence :} 4-17 secondes par requête
    \item \textbf{Précision :} $\sim$85\% pour les requêtes complexes
    \item \textbf{Fallback :} Les providers de secours ne sont pas configurés
\end{enumerate}
\end{warningbox}

\newpage

%==============================================================================
% SECTION 8 : RECOMMANDATIONS FUTURES
%==============================================================================
\section{\faRocket\ Recommandations Futures}

\subsection{Phase 1 : Court Terme - Amélioration du LLM}

\begin{table}[H]
\centering
\rowcolors{2}{lightgray}{white}
\begin{tabular}{|>{\bfseries}l|c|c|c|c|}
\hline
\rowcolor{primaryblue}
\textcolor{white}{\textbf{Option}} & \textcolor{white}{\textbf{Coût/mois}} & \textcolor{white}{\textbf{Latence}} & \textcolor{white}{\textbf{Précision}} & \textcolor{white}{\textbf{Recommandé}} \\
\hline
\cellcolor{successbox!20}Gratuit (Actuel) & 0€ & 4-17s & $\sim$85\% & Dev \\
\hline
\rowcolor{successbox!30}
\cellcolor{successbox!40}OpenAI GPT-4o-mini & \textbf{15-20€} & \textbf{1-3s} & \textbf{$\sim$95\%} & \faCheckCircle\ \textbf{Production} \\
\hline
Claude 3.5 Sonnet & 40-50€ & 2-4s & $\sim$98\% & Premium \\
\hline
Azure OpenAI & 20-30€ & 1-2s & $\sim$95\% & Enterprise \\
\hline
\end{tabular}
\caption{Options de LLM pour Production}
\end{table}

\begin{successbox}[title={\faStar\ Recommandation}]
\textbf{OpenAI GPT-4o-mini} est le meilleur compromis coût/performance pour une mise en production.
\begin{itemize}
    \item Budget : $\sim$15-20€/mois pour 500 requêtes/jour
    \item Latence réduite de 80\% (1-3s vs 4-17s)
    \item Précision améliorée de 10\% ($\sim$95\% vs $\sim$85\%)
\end{itemize}
\end{successbox}

\subsection{Phase 2 : Moyen Terme - Optimisations}

\begin{enumerate}[leftmargin=*]
    \item \textbf{\faMemory\ Cache Redis}
    \begin{itemize}
        \item Cacher les réponses fréquentes
        \item Réduction de 80\% des appels LLM
        \item Temps de réponse < 100ms pour requêtes cachées
    \end{itemize}
    
    \item \textbf{\faChartBar\ Monitoring et Analytics}
    \begin{itemize}
        \item Dashboard de suivi des requêtes
        \item Alertes sur les erreurs
        \item Métriques de performance
    \end{itemize}
    
    \item \textbf{\faGraduationCap\ Fine-tuning}
    \begin{itemize}
        \item Entraîner un modèle spécialisé GMAO
        \item Améliorer la précision des requêtes SQL
        \item Réduire les coûts à long terme
    \end{itemize}
\end{enumerate}

\subsection{Phase 3 : Long Terme - Production}

\begin{enumerate}[leftmargin=*]
    \item \textbf{\faCloud\ Infrastructure Cloud}
    \begin{itemize}
        \item Déploiement sur Azure/AWS
        \item Load balancing et auto-scaling
        \item Haute disponibilité (99.9\%)
    \end{itemize}
    
    \item \textbf{\faShieldAlt\ Sécurité Renforcée}
    \begin{itemize}
        \item Audit complet des requêtes
        \item Chiffrement end-to-end
        \item Conformité RGPD
    \end{itemize}
    
    \item \textbf{\faChartArea\ Fonctionnalités Avancées}
    \begin{itemize}
        \item Graphiques et visualisations
        \item Export PDF/Excel
        \item Alertes proactives
    \end{itemize}
\end{enumerate}

\newpage

%==============================================================================
% SECTION 9 : ESTIMATION DES COÛTS
%==============================================================================
\section{\faEuroSign\ Estimation des Coûts}

\subsection{Scénario : 500 requêtes/jour}

\begin{figure}[H]
\centering
\begin{tikzpicture}
    \begin{axis}[
        ybar,
        width=14cm,
        height=8cm,
        ylabel={Coût mensuel (€)},
        symbolic x coords={Gratuit, GPT-4o-mini, Claude 3.5, Azure OpenAI},
        xtick=data,
        nodes near coords,
        nodes near coords align={vertical},
        ymin=0,
        ymax=50,
        bar width=1.5cm,
        enlarge x limits=0.2,
        legend style={at={(0.5,-0.15)}, anchor=north, legend columns=-1},
    ]
    \addplot[fill=successbox!70] coordinates {(Gratuit,0) (GPT-4o-mini,17) (Claude 3.5,45) (Azure OpenAI,25)};
    \end{axis}
\end{tikzpicture}
\caption{Comparaison des Coûts Mensuels par Provider}
\end{figure}

\begin{table}[H]
\centering
\rowcolors{2}{lightgray}{white}
\begin{tabular}{|l|c|c|c|c|}
\hline
\rowcolor{primaryblue}
\textcolor{white}{\textbf{Configuration}} & \textcolor{white}{\textbf{Coût/mois}} & \textcolor{white}{\textbf{Latence}} & \textcolor{white}{\textbf{Précision}} & \textcolor{white}{\textbf{Usage}} \\
\hline
Gratuit (Actuel) & 0€ & 4-17s & $\sim$85\% & Développement \\
\rowcolor{successbox!20}
OpenAI GPT-4o-mini & $\sim$15€ & 1-3s & $\sim$95\% & \textbf{Production} \\
Claude 3.5 Sonnet & $\sim$40€ & 2-4s & $\sim$98\% & Premium \\
Azure OpenAI & $\sim$20€ & 1-2s & $\sim$95\% & Enterprise \\
\hline
\end{tabular}
\caption{Résumé des Options}
\end{table}

%==============================================================================
% CONCLUSION
%==============================================================================
\section{\faFlagCheckered\ Conclusion}

\begin{tcolorbox}[
    enhanced,
    colback=primaryblue!10,
    colframe=primaryblue,
    title={\faStar\ Points Clés à Retenir},
    fonttitle=\bfseries
]
\begin{enumerate}
    \item \textbf{Architecture Text-to-SQL} choisie car idéale pour les données structurées GMAO
    \item \textbf{Pas de RAG} car trop coûteux et moins précis pour les requêtes SQL
    \item \textbf{Approche hybride} à 3 niveaux pour optimiser coûts et performances
    \item \textbf{Configuration actuelle gratuite} adaptée au développement
    \item \textbf{Recommandation production} : OpenAI GPT-4o-mini ($\sim$15-20€/mois)
\end{enumerate}
\end{tcolorbox}

\vspace{1cm}

\begin{center}
\begin{tikzpicture}
    \node[draw=secondarygreen, fill=secondarygreen!20, rounded corners=10pt,
          minimum width=10cm, minimum height=1.5cm, align=center] {
        {\Large\bfseries\color{secondarygreen} Le GMAO Chatbot est opérationnel !}\\
        Prêt pour la démonstration et les tests.
    };
\end{tikzpicture}
\end{center}

%==============================================================================
% FIN DU DOCUMENT
%==============================================================================
\end{document}
